\documentclass[a4paper,11pt]{article}
\usepackage[utf8]{inputenc}
\usepackage[T1]{fontenc}
\usepackage[french]{babel}
\usepackage[left=1cm,right=1cm,top=.5cm,bottom=0cm]{geometry}
\usepackage{enumitem}
\usepackage{hyperref}
\usepackage{xcolor}
\usepackage{titlesec}
\usepackage{fontawesome5}
\usepackage{lmodern}
\usepackage[normalem]{ulem}

% --- Couleurs Cotral Lab ---
\definecolor{primary}{RGB}{0, 80, 140}
\definecolor{secondary}{RGB}{80, 80, 80}

\hypersetup{colorlinks=true, linkcolor=primary, filecolor=primary, urlcolor=primary}
\titleformat{\section}{\normalsize\bfseries\color{primary}\uppercase}{}{0em}{}[\titlerule]
\titlespacing{\section}{0pt}{6pt}{3pt}

\begin{document}

% --- EN-TÊTE ---
\begin{center}
    {\large \textbf{NJIHA KUITO Brayanne Stella}} \\ \vspace{2pt}
    \textbf{Alternante Ingénieure R\&D -- Electronique Analogique \& Systèmes Embarqués} \\
    \textit{\small Disponible dès septembre 2026 -- Alternance 1 an} \\ \vspace{2pt}
    \small
    \faMapMarker\ Cergy (95), France \quad
    \faPhone\ (+33) 6 52 31 08 46 \quad
    \faEnvelope\ \href{mailto:stellanjiha1@gmail.com}{stellanjiha1@gmail.com} \quad
    \faCar\ Permis B
\end{center}

\vspace{-2pt}

% --- PROFIL ---
\noindent\small{Etudiante en électronique avec une expérience en conception PCB (KiCad), simulation de circuits analogiques (LTSpice/PSpice), développement embarqué (C, STM32) et validation fonctionnelle en laboratoire. Curieuse et rigoureuse, de la schématique jusqu'aux essais de validation.}

\vspace{2pt}

% --- FORMATION ---
\section{Formation}
\begin{itemize}[leftmargin=*, label=\tiny$\bullet$, nosep]
    \item \textbf{CY Paris Université} \hfill \textit{2026 -- 2027} \\
    \small{Master 2 EEA -- Electronique, Energie Electrique, Automatique.}
    \item \textbf{ENSEA -- Ecole Nationale Supérieure de l'Electronique et de ses Applications} \hfill \textit{2025 -- 2026} \\
    \small{Master 1 Electronique (échange) -- Electronique de puissance, Conversion d'énergie, CEM, Véhicules électriques.}
    \item \textbf{Ecole Nationale Supérieure Polytechnique de Yaoundé (ENSPY)} \hfill \textit{2022 -- 2025} \\
    \small{Cursus Ingénieur Génie Electrique -- Electrotechnique, Energies renouvelables, Electronique analogique et numérique, Automatique.}
\end{itemize}

\vspace{1pt}

% --- COMPÉTENCES ---
\section{Compétences Techniques}
\begin{itemize}[leftmargin=*, nosep]
    \item \small{\textbf{Conception \& CAO :} KiCad (schématique, routage PCB, dossiers de définition), électronique analogique faible puissance, convertisseurs DC/DC (Buck, Boost, Flyback), onduleurs, choix et dimensionnement de composants.}
    \item \small{\textbf{Simulation :} LTSpice, PSpice -- modélisation dynamique, stabilité, rendement, pertes MOSFET/IGBT, études de faisabilité.}
    \item \small{\textbf{Validation \& Instrumentation :} Oscilloscope, GBF, multimètre, analyseur de signaux, plans de test, analyse des écarts, rapports de conformité.}
    \item \small{\textbf{Embarqué \& Programmation :} C/C++, VHDL, STM32, Arduino, Raspberry Pi, I2C/UART, debug et mise au point prototype.}
    \item \small{\textbf{CEM :} Couplages EM, plans de masse, immunité des signaux. \quad \textbf{Langues :} Français (maternel), Anglais (professionnel).}
\end{itemize}

\vspace{1pt}

% --- PROJETS ---
\section{Projets Académiques}
\begin{itemize}[leftmargin=*, label={-}, itemsep=1pt]

\item \textbf{Robot à Navigation Autonome -- Conception PCB \& Embarqué} \hfill \textit{2025 -- 2026} \\
\textit{ENSEA} \hfill \textit{STM32, Raspberry Pi 4, KiCad, I2C, C, ROS2, LIDAR}
\vspace{-2pt}
    \begin{itemize}[leftmargin=0.4cm, nosep, label=\tiny$\bullet$]
        \item \small{Conception d'une architecture embarquée STM32/Raspberry Pi 4 via bus I2C/UART, avec intégration de capteurs (LIDAR, encodeurs) et actionneurs pour la navigation autonome.}
        \item \small{Saisie du schéma de principe et routage complet sous KiCad d'une carte d'alimentation (régulateur de tension), en tenant compte des contraintes CEM et de puissance.}
        \item \small{Assemblage, soudure de précision, debug et validation fonctionnelle en laboratoire. Programmation des lois de commande moteur en C sur STM32.}
    \end{itemize}

\item \textbf{Simulation Convertisseurs DC/DC \& Onduleurs} \hfill \textit{2025 -- 2026} \\
\textit{ENSEA} \hfill \textit{LTSpice, PSpice}
\vspace{-2pt}
    \begin{itemize}[leftmargin=0.4cm, nosep, label=\tiny$\bullet$]
        \item \small{Modélisation structurelle et simulation dynamique de topologies analogiques : Buck, Boost et onduleurs monophasés.}
        \item \small{Calculs de rendement, analyse de stabilité des signaux, estimation des pertes par commutation/conduction (MOSFET/IGBT) et rédaction de rapports de faisabilité.}
    \end{itemize}

\item \textbf{Caractérisation CEM -- Lignes Microrubans \& Plans de Masse} \hfill \textit{2025 -- 2026} \\
\textit{ENSEA} \hfill \textit{Oscilloscope, GBF, Analyseur de signaux}
\vspace{-2pt}
    \begin{itemize}[leftmargin=0.4cm, nosep, label=\tiny$\bullet$]
        \item \small{Etude expérimentale des couplages EM sur lignes microrubans et analyse de l'impact des plans de masse sur l'immunité des signaux de cartes électroniques.}
        \item \small{Définition d'un plan de mesures, validation instrumentale, analyse des écarts et rédaction de rapports de conformité.}
    \end{itemize}

\end{itemize}

\vspace{1pt}

% --- EXPÉRIENCES ---
\section{Expériences Professionnelles}
\begin{itemize}[leftmargin=*, label={-}, itemsep=1pt]

\item \textbf{Stagiaire Systèmes Embarqués -- ENSPY} \hfill \textit{Juin -- Août 2025}
\vspace{-2pt}
    \begin{itemize}[leftmargin=0.4cm, nosep, label=\tiny$\bullet$]
        \item \small{Développement en C sur Arduino pour véhicule autonome : intégration capteurs/actionneurs, optimisation des algorithmes de détection et contrôle moteur.}
        \item \small{Tests fonctionnels et validation de prototype, rédaction de la documentation technique et traçabilité des essais.}
    \end{itemize}

\item \textbf{Stagiaire Ouvrier -- OMNIUM SERVICES}
\vspace{-2pt}
    \begin{itemize}[leftmargin=0.4cm, nosep, label=\tiny$\bullet$]
        \item \small{Repérage, étiquetage et raccordement de câbles électriques, assistance aux tests d'équipements et participation aux vérifications de conformité.}
    \end{itemize}

\end{itemize}

\vspace{1pt}

% --- ACTIVITÉS ---
\section{Activités \& Engagements}
\begin{itemize}[leftmargin=*, nosep, label={-}]
    \item \small{\textbf{Vice-Présidente -- Club Génie Electrique | ENSPY} : Coordination du bureau, organisation d'événements techniques, représentation officielle, développement du leadership et de la prise d'initiative.}
    \item \small{\textbf{Cheffe de la Cellule SocioCulturelle -- AFENSPY} : Organisation logistique de grands événements culturels, coordination d'équipes et planification.}
\end{itemize}

\vspace{1pt}

% --- ATOUTS & CENTRES D'INTÉRÊT ---
\section{Atouts \& Centres d'intérêt}
\begin{itemize}[leftmargin=*, nosep, label=\tiny$\bullet$]
    \item \small{\textbf{Atouts :} Curiosité technique, rigueur, esprit d'analyse, autonomie, force de proposition, travail en équipe, goût pour la résolution de problèmes complexes.}
    \item \small{\textbf{Centres d'intérêt :} Escalade \quad|\quad Handball \quad|\quad Musique \quad|\quad Danse \quad|\quad Mangas}
\end{itemize}

\end{document}
